\documentclass[11pt,letterpaper]{article}

% --- Packages ---
\usepackage[margin=1in]{geometry}
\usepackage{booktabs}
\usepackage[colorlinks=true,linkcolor=blue,urlcolor=blue,citecolor=blue]{hyperref}
\usepackage{amsmath}
\usepackage{listings}
\usepackage{xcolor}
\usepackage{graphicx}
\usepackage{siunitx}
\usepackage{enumitem}
\usepackage{longtable}
\usepackage{caption}

% --- Listings style ---
\lstset{
  basicstyle=\ttfamily\small,
  backgroundcolor=\color{gray!10},
  frame=single,
  framerule=0pt,
  breaklines=true,
  columns=fullflexible,
  keepspaces=true,
  aboveskip=1em,
  belowskip=1em,
}

% --- siunitx setup ---
\sisetup{
  per-mode=symbol,
  detect-all,
}

% --- Title ---
\title{Summer Measurement Plan: Concrete Goals for 2026}
\author{LDRD FFA@CEBAF Permanent Magnet Radiation Damage Study}
\date{June 2026}

\begin{document}
\maketitle

\noindent\textbf{Interns:} Moishe, Evelyn \\
\textbf{Duration:} ${\sim}$10 weeks (June--August 2026)

\bigskip\noindent\textbf{Figures:} This document references PNG files from the
project. For Overleaf, upload them to the same directory as this \texttt{.tex}
file. Source locations relative to \texttt{Claude\_Analysis1/}:
\begin{itemize}[nosep]
  \item \texttt{2026\_Data\_Run/Analysis/} (C2-series, JA-series)
\end{itemize}

\tableofcontents
\newpage

%=============================================================================
\section{Overview}
\label{sec:overview}
%=============================================================================

This plan has two complementary goals:

\begin{enumerate}
  \item \textbf{Improve measurement reliability:} Accumulate multi-session
    baselines, characterize temperature sensors, and quantify noise floors for
    both Helmholtz and Teslameter instruments.

  \item \textbf{Enable dual-measurement publications:} Get the Teslameter
    precise enough to serve as an independent primary measurement, not just a
    cross-check. If we can show demagnetization in both integrated moment
    (Helmholtz) and spatially-resolved point field (Teslameter), across all
    three sample types (Y-plates, H-plates, A-samples), we are positioned for
    high-impact journal publications. H-plates are especially important because
    they directly represent Halbach/FFA magnet geometry.
\end{enumerate}

The plan is organized into 6 priorities. All priorities can overlap; the ordering
reflects when each should start.

\subsection{Equipment Available}

\begin{itemize}[nosep]
  \item Helmholtz coil + Fluxmeter Model 2130 (integrated magnetic moment, mWC)
  \item SENIS 3MH6-E Teslameter + Hall Probe Type C (3-axis field, mT)
  \item Arduino AM2302 temperature/humidity sensor (ambient temp + RH)
  \item 3D-printed Teslameter sample rigs (top, front, side faces)
  \item Lab Y-plates: Y-\{8,\,14,\,27,\,28,\,29,\,31,\,33,\,35,\,37\} (9
    plates, 36 slots)
  \item Lab H-plates: 24 north + 24 south (48 plates, 192 slots)
  \item Lab A-samples: all lab A-sample pairs
\end{itemize}

\subsection{General Protocol (Every Session)}

\begin{enumerate}[nosep]
  \item Power on Helmholtz and Teslameter. Allow 15+ minutes for thermal
    equilibration.
  \item Start Arduino sensor logging.
  \item Measure Y-14 calibration plate (all 4 slots, Helmholtz + Teslameter top).
  \item Measure all scheduled samples.
  \item Measure Y-14 again at end of session.
  \item Record session notes: start/end times, any anomalies, room conditions.
  \item Verify all data files saved and named correctly before shutdown.
\end{enumerate}

Figure~\ref{fig:c2-cal} shows the Y-14 repeatability achieved in April 2026.

\begin{figure}[htbp]
  \centering
  \includegraphics[width=0.85\textwidth]{C2-1_calibration_repeatability.png}
  \caption{Y-14 calibration plate repeatability from the April 2026 measurement
    campaign. Corrected std per slot: 0.06--0.27\%.}
  \label{fig:c2-cal}
\end{figure}


%=============================================================================
\section{Priority 1: Lab Control Baseline Accumulation (Weeks 1--5)}
\label{sec:p1}
%=============================================================================

\subsection{Goal}

Measure every lab plate in 5--10 independent sessions to build a high-statistics
baseline that reduces per-plate noise as $1/\sqrt{N}$.

\subsection{What to Measure}

\begin{table}[htbp]
\centering
\caption{Per-session measurement counts for Priority~1.}
\label{tab:p1-counts}
\begin{tabular}{@{}lccc@{}}
\toprule
Sample type & Count & Slots per plate & Total measurements per session \\
\midrule
Lab Y-plates  & 9  & 4 & 36 Helmholtz + 36 Teslameter (top) \\
Lab H-plates  & 48 & 4 & 192 Helmholtz + 192 Teslameter (top) \\
Lab A-samples & All pairs & 2 per pair & ${\sim}$96 Helmholtz + ${\sim}$96 Teslameter \\
\bottomrule
\end{tabular}
\end{table}

\textbf{Instruments per sample:} Helmholtz (mWC) + Teslameter top face + Arduino
temp.

\subsection{Cadence}

${\sim}$2 sessions per week for 5 weeks = 10 sessions.

Each session takes approximately:
\begin{itemize}[nosep]
  \item Y-plates: 9 plates $\times$ 4 slots $\times$ ${\sim}$2 min = ${\sim}$72 min
  \item H-plates: 48 plates $\times$ 4 slots $\times$ ${\sim}$1.5 min = ${\sim}$288 min (${\sim}$5 hours)
  \item A-samples: ${\sim}$96 pairs $\times$ ${\sim}$1.5 min = ${\sim}$144 min (${\sim}$2.5 hours)
  \item Y-14 calibration: 2 $\times$ ${\sim}$8 min = 16 min
  \item Total: ${\sim}$8--9 hours per full session
\end{itemize}

\textbf{Recommendation:} Split into half-sessions if needed. Y-plates + Y-14 cal
can be a morning session (${\sim}$1.5 hours). H-plates can be an afternoon
session. A-samples on alternating days.

\textbf{Important:} Do not neglect H-plates and A-samples in favor of Y-plates.
For the broader publication goals, H-plate and A-sample baselines are equally
critical. H-plates demonstrate degradation in Halbach (FFA-relevant) geometry;
A-samples provide paired field-vs-zero-field comparisons. Both need the same
multi-session baseline accumulation as Y-plates. The weekly schedule should
ensure all three sample types get equal attention.

\subsection{Statistical Target}

With $N=10$ independent sessions per lab Y-plate:

\begin{table}[htbp]
\centering
\caption{Statistical improvement from multi-session baselines.}
\label{tab:p1-stats}
\begin{tabular}{@{}lccc@{}}
\toprule
Metric & Current ($N \sim 1$--$2$) & Target ($N=10$) & Improvement \\
\midrule
Per-plate baseline std & ${\sim}$0.2--0.3\% & ${\sim}$0.2--0.3\% (unchanged) & -- \\
Per-plate baseline SEM & ${\sim}$0.2--0.3\% & ${\sim}$0.06--0.10\%            & 3--5$\times$ \\
Lab fleet differential SEM & 0.038\%        & ${\sim}$0.012\%                 & ${\sim}3\times$ \\
\bottomrule
\end{tabular}
\end{table}

The lab null becomes definitive: if the lab differential is truly zero, we will
measure $0.000\% \pm 0.012\%$, ruling out any hidden systematic $>$0.024\%
($2\sigma$).

\subsection{Data Naming Convention}

Follow existing convention: \texttt{\{sample\_id\}\_helmholtzTemp.dat} for
Arduino data, \texttt{\{sample\_id\}.dat} for Helmholtz,
\texttt{\{sample\_id\}\_\{face\}.dat} for Teslameter.

Organize by date: \texttt{2026\_Data\_Run/YYYY-MM-DD\_Helmholtz/} and
\texttt{2026\_Data\_Run/YYYY-MM-DD\_Teslameter/}.


%=============================================================================
\section{Priority 2: Temperature Cross-Calibration (Weeks 2--5)}
\label{sec:p2}
%=============================================================================

\subsection{Goal}

Characterize the systematic difference between the Arduino AM2302 ambient sensor
and the Teslameter probe temperature, as a function of temperature, time of day,
and humidity.

\subsection{Protocol}

Measure Y-14 (calibration plate, all 4 slots) at multiple times throughout the
day for 5+ separate days:

\begin{table}[htbp]
\centering
\caption{Temperature cross-calibration measurement schedule.}
\label{tab:p2-schedule}
\begin{tabular}{@{}lp{9cm}@{}}
\toprule
Time & Purpose \\
\midrule
Morning (07:00--08:00)     & Cold start; largest expected sensor disagreement (sample may lag air) \\
Mid-morning (09:00--10:00) & Room warming up \\
Midday (12:00--13:00)      & Near steady-state \\
Afternoon (15:00--16:00)   & Warmest; sensors likely most equilibrated \\
\bottomrule
\end{tabular}
\end{table}

For each measurement:
\begin{enumerate}[nosep]
  \item Record Arduino temperature and humidity (from \texttt{*\_helmholtzTemp.dat})
  \item Record Teslameter temperature (from \texttt{*\_top.dat})
  \item Note time since power-on of instruments
  \item Note time since samples were placed on bench (equilibration time)
\end{enumerate}

\subsection{Days to Target}

\begin{itemize}[nosep]
  \item At least 2 ``cold start'' days: arrive early, measure before room warms up
  \item At least 2 ``hot'' days: ambient $>$\SI{27}{\degreeCelsius} (expected in July--August)
  \item At least 1 day with humidity $>$55\% (if possible)
  \item At least 5 total days with 3+ measurements per day
\end{itemize}

\subsection{Deliverables}

\begin{enumerate}[nosep]
  \item Characterization of the Arduino--Teslameter offset: is it constant,
    linear in temperature, or more complex? (Do NOT assume either sensor is
    ground truth.)
  \item Assessment of which sensor (or combination) better predicts Helmholtz
    residuals after temperature correction.
  \item Humidity effect: does humidity correlate with temperature disagreement?
  \item Equilibration study: how long after placing samples on the bench do the
    two sensors converge? This constrains whether the offset is due to sample
    thermal lag vs sensor bias.
  \item Draft recommendation: what temperature source should be used for the
    correction formula, and with what estimated uncertainty?
\end{enumerate}

\textbf{Key caution:} The Arduino measures ambient air temperature, which can be
influenced by cross-breezes, laptop fans, HVAC airflow, and distance from the
sample. The Teslameter probe is in contact with the magnet surface, so it may be
closer to the actual sample temperature in some conditions. Do not assume either
is ``correct'' without evidence.

\subsection{Statistical Target}

Understand the offset well enough to reduce temperature uncertainty to
$<$\SI{0.3}{\degreeCelsius}, which would bring the temperature systematic on the
differential from $\pm 0.033\%$ to $\pm 0.020\%$.


%=============================================================================
\section{Priority 3: Teslameter Repeatability Study (Weeks 3--7)}
\label{sec:p3}
%=============================================================================

\subsection{Goal}

Quantify the Teslameter positioning noise per face, per material, per sample,
using a remove-and-reinsert protocol.

\subsection{Sample Selection}

8 representative samples spanning all three sample types and all four materials:

\begin{table}[htbp]
\centering
\caption{Samples for Teslameter repeatability study.}
\label{tab:p3-samples}
\begin{tabular}{@{}clllp{5cm}@{}}
\toprule
\# & Sample & Type & Material & Why \\
\midrule
1 & Y-14-1 & Y-plate       & (check assignment) & Calibration plate, well-characterized \\
2 & Y-14-2 & Y-plate       & (check assignment) & Same plate, different material \\
3 & Y-14-3 & Y-plate       & (check assignment) & Same plate, different material \\
4 & Y-14-4 & Y-plate       & (check assignment) & Same plate, different material \\
5 & One lab Hn & H-plate (north) & Mixed & Different rig geometry from Y-plates \\
6 & One lab Hs & H-plate (south) & Mixed & North vs south comparison \\
7 & One A-sample pair & A-sample & NdFeB & Smallest sample, potentially worst positioning \\
8 & One A-sample pair & A-sample & SmCo & Material comparison for A-samples \\
\bottomrule
\end{tabular}
\end{table}

Including H-plate and A-sample geometries is essential because the rig
positioning noise may differ by sample type (different rig, different magnet
shape, different field gradients). If we only characterize Y-plate noise, we
cannot make claims about Teslameter precision on H-plates or A-samples.

\subsection{Protocol (Per Sample, Per Day)}

\begin{enumerate}[nosep]
  \item Insert sample into rig for top face measurement
  \item Record Teslameter reading ($B_x$, $B_y$, $B_z$, $T$)
  \item Remove sample completely from rig
  \item Re-insert sample
  \item Record reading again
  \item Repeat steps 3--5 until you have 10 readings per face
  \item Move to front face rig; repeat 10$\times$
  \item Move to side face rig; repeat 10$\times$
  \item Total: 30 readings per sample per day
\end{enumerate}

\subsection{Schedule}

Run on 3 different days (ideally different weeks) to test day-to-day
reproducibility.

Total: 8 samples $\times$ 3 faces $\times$ 10 insertions $\times$ 3 days = 720
readings. Time estimate: ${\sim}$10 min per sample per face per day, ${\sim}$4
hours per day, ${\sim}$12 hours total across 3 days.

\subsection{Analysis}

For each sample/face combination:
\begin{itemize}[nosep]
  \item Mean, std, CV (coefficient of variation $= \mathrm{std}/\mathrm{mean} \times 100\%$)
  \item Within-day vs between-day variance (is day-to-day variation larger?)
  \item Compare to existing estimates (top 0.14--0.29\%, front 0.57--0.91\%, side
    0.90--1.90\%)
\end{itemize}

\subsection{Decision Point}

If top-face CV $< 0.15\%$: current rig is adequate for top-face trend tracking.
If top-face CV $> 0.30\%$: pursue machined jig (Priority~6 stretch goal).


%=============================================================================
\section{Priority 4: Helmholtz Repeatability Study (Weeks 5--8)}
\label{sec:p4}
%=============================================================================

\subsection{Goal}

Quantify the single-insertion reproducibility of the Helmholtz coil measurement.

\subsection{Samples}

\begin{itemize}[nosep]
  \item Y-14 (calibration plate, all 4 slots)
  \item 2 additional lab Y-plates (e.g., Y-8, Y-27): 8 more slots
  \item Total: 12 slots
\end{itemize}

\subsection{Protocol (Per Session)}

\begin{enumerate}[nosep]
  \item Measure all 12 slots normally (insert, record mWC, remove)
  \item Re-measure all 12 slots (insert, record mWC, remove)
  \item Repeat until each slot has 10 readings in the same session
  \item Record Arduino temperature throughout (should be stable if room is
    equilibrated)
  \item Total: 120 Helmholtz readings in a single session
\end{enumerate}

\subsection{Time Estimate}

${\sim}$1 minute per insertion $\times$ 120 insertions = ${\sim}$2 hours per
session. Do 2--3 sessions on different days.

\subsection{Analysis}

For each slot:
\begin{itemize}[nosep]
  \item 10-reading mean, std, CV
  \item Compare to nominal instrument accuracy (0.05\%)
  \item Check for systematic drift within the 10-reading sequence (does the
    reading drift up or down as the session progresses?)
\end{itemize}

\subsection{Expected Result}

Single-insertion std of 0.02--0.05\%. This would confirm that Helmholtz
positioning noise is negligible compared to gain drift (0.5--0.7\%) and
statistical scatter (0.17--0.20\%).


%=============================================================================
\section{Priority 5: Environmental Monitoring (Weeks 6--9)}
\label{sec:p5}
%=============================================================================

\subsection{Goal}

Characterize how room environment (temperature, humidity) varies during
measurement sessions and how it correlates with Helmholtz readings.

\subsection{Protocol}

\begin{enumerate}[nosep]
  \item Arduino sensor logs continuously during all measurement sessions (it
    should already be doing this for Priority~1)
  \item Additionally, run Arduino logging on 3--5 non-measurement days to capture
    baseline environmental patterns
  \item Record HVAC status (on/off, set point) if accessible
\end{enumerate}

\subsection{Analysis}

\begin{itemize}[nosep]
  \item Temperature profile: typical daily range, rate of change, time-of-day
    pattern
  \item Humidity profile: range, correlation with temperature
  \item Helmholtz drift within sessions vs temperature change
  \item Is there a measurable correlation between Helmholtz reading and Arduino
    temperature beyond what the temperature correction formula predicts?
\end{itemize}

Figure~\ref{fig:c2-temp} shows the temperature swing from the April 2026 campaign
as an example of the variability that must be characterized.

\subsection{Deliverable}

Environmental characterization report with plots:
\begin{itemize}[nosep]
  \item Temperature vs time-of-day (5+ days overlaid)
  \item Humidity vs time-of-day
  \item Helmholtz residual (after temp correction) vs session temperature change
  \item Thermal equilibration timescale for samples (from Priority~2 data)
\end{itemize}

\begin{figure}[htbp]
  \centering
  \includegraphics[width=0.85\textwidth]{C2-2_temperature_evolution.png}
  \caption{Temperature evolution during the April 2026 measurement campaign
    (Campaign~2). An \SI{8.6}{\degreeCelsius} swing was observed over the
    measurement day.}
  \label{fig:c2-temp}
\end{figure}


%=============================================================================
\section{Priority 6: Improved-Protocol Re-measurement (Weeks 8--10)}
\label{sec:p6}
%=============================================================================

\subsection{Goal}

If Priorities 1--4 are complete, re-measure all 30 tunnel Y-plates using the
improved protocol (multi-insertion Helmholtz, dual temperature sensors, Y-14
calibration at start and end).

\subsection{Protocol}

For each tunnel Y-plate (4 slots):
\begin{enumerate}[nosep]
  \item Helmholtz: 3 insertions per slot, record all 3 readings
  \item Arduino temperature recorded for each insertion
  \item Teslameter top face: 3 insertions per slot
  \item Y-14 measured at start and end of every batch of 5 plates
\end{enumerate}

\subsection{Deliverable}

``C2 improved'' dataset: 120 tunnel Y-plate measurements with:
\begin{itemize}[nosep]
  \item Multi-insertion Helmholtz means and stds
  \item Dual temperature readings
  \item Y-14 calibration brackets
  \item Full Arduino environmental log
\end{itemize}

\subsection{Time Estimate}

30 plates $\times$ 4 slots $\times$ 3 insertions $\times$ ${\sim}$2 min =
${\sim}$12 hours of Helmholtz. Plus Teslameter: ${\sim}$12 hours. Plus
calibration: ${\sim}$3 hours. Total: ${\sim}$27 hours, spread over ${\sim}$4
days.


%=============================================================================
\section{Stretch Goals (If Time Permits)}
\label{sec:stretch}
%=============================================================================

\subsection{S1: Prototype Machined Teslameter Jig}

If Priority~3 shows top-face CV $> 0.30\%$, design and fabricate an aluminum jig
for top-face measurements. Requirements:
\begin{itemize}[nosep]
  \item Repeatable sample positioning to $<$\SI{0.05}{\milli\meter}
  \item Easy insertion/removal (no tools required during measurement)
  \item Compatible with Y-plate and H-plate sample dimensions
\end{itemize}

This is the single highest-leverage stretch goal: if the rig is the noise
bottleneck, fixing it unlocks the Teslameter as a primary science measurement.

\subsection{S2: Humidity--Field Investigation}

With Arduino humidity data: is there a measurable correlation between relative
humidity and Helmholtz or Teslameter readings? Anecdotally, magnetic measurements
can be affected by condensation or air moisture, but this has never been tested
on our setup.

\subsection{S3: Tunnel H-Plate and A-Sample Re-measurement}

Re-measure all tunnel H-plates and A-samples with the improved protocol
(multi-insertion, dual temp sensors). Combined with the lab baselines from
Priority~1 and the Teslameter noise characterization from Priority~3, this
creates a dataset suitable for dual-measurement (point + integrated) publication
on H-plate and A-sample degradation.


%=============================================================================
\section{Weekly Schedule Template}
\label{sec:schedule}
%=============================================================================

\begin{longtable}{@{}cp{2.8cm}p{2.4cm}p{2.4cm}p{2.8cm}@{}}
\caption{Weekly schedule template. P1 = Baselines, P2 = Temp Cal, P3 = Teslameter
  Repeatability.} \label{tab:schedule} \\
\toprule
Week & P1 (Baselines) & P2 (Temp Cal) & P3 (Tesla Rep.) & Other \\
\midrule
\endfirsthead
\toprule
Week & P1 (Baselines) & P2 (Temp Cal) & P3 (Tesla Rep.) & Other \\
\midrule
\endhead
1  & Sessions 1--2 (Y + Y-14 cal; learn protocol) & -- & -- & Training, setup \\
2  & Sessions 3--4 (Y + H-plates) & Start temp cal & -- & -- \\
3  & Sessions 5--6 (Y + H + A, full) & Continue temp cal & Start repeatability & -- \\
4  & Sessions 7--8 (Y + H + A, full) & Wrap up temp cal & Continue repeatability & -- \\
5  & Sessions 9--10 (Y + H + A, full) & Final day if needed & Continue repeatability & Start P4 pilot \\
6  & Catch-up if needed & -- & Wrap up & Helmholtz repeatability \\
7  & Additional H/A sessions if behind & -- & Analysis + writeup & Helmholtz repeatability \\
8  & -- & -- & -- & Begin improved re-measurement \\
9  & Wrap-up & -- & -- & Improved re-measurement \\
10 & -- & -- & -- & Final analysis, reports \\
\bottomrule
\end{longtable}

Note: Weeks 1--2 ramp up gradually (Y-plates first to learn the protocol, then
add H-plates and A-samples). By week~3, every session should include all three
sample types. Do not let H-plates and A-samples fall behind Y-plates.


%=============================================================================
\section{Success Metrics}
\label{sec:success}
%=============================================================================

By the end of summer, we should have:

\begin{enumerate}
  \item \textbf{10 independent measurement sessions} for all 9 lab Y-plates, all
    48 lab H-plates, and lab A-samples, producing a lab Y-plate differential SEM
    of ${\sim}$0.012\% and meaningful H-plate/A-sample baselines

  \item \textbf{Temperature cross-calibration} characterizing the
    Arduino--Teslameter offset with a clear recommendation for which sensor (or
    combination) to use, targeting temperature uncertainty
    $<$\SI{0.3}{\degreeCelsius} ($\pm 0.020\%$ on differential)

  \item \textbf{Teslameter positioning noise} quantified per face, per material,
    per sample type (Y, H, A), from a 720-reading dataset across 3 days.
    Determines whether the Teslameter can serve as an independent primary
    measurement for publication.

  \item \textbf{Helmholtz single-insertion std} measured (expect 0.02--0.05\%),
    confirming negligible positioning noise

  \item \textbf{Environmental characterization} showing typical
    temperature/humidity patterns during measurement sessions

  \item (Stretch) \textbf{Improved-protocol tunnel Y-plate dataset} for C2
    comparison
\end{enumerate}

\subsection{Publication Readiness Checkpoint}

At the end of the summer, evaluate:
\begin{itemize}[nosep]
  \item Can we show consistent demagnetization in both Helmholtz and Teslameter
    data on Y-plates? (Requires top-face Teslameter noise $<$0.15\%)
  \item Do H-plate results, measured with improved baselines and both
    instruments, support the Y-plate story? (Requires multi-session H-plate
    baselines)
  \item Are A-sample paired comparisons informative? (Requires A-sample
    baselines)
\end{itemize}

If all three are ``yes,'' we have the dataset for a high-impact journal paper
showing dual-measurement, multi-geometry evidence for radiation-induced permanent
magnet demagnetization in an operating accelerator.


%=============================================================================
\section{Important Reminders}
\label{sec:reminders}
%=============================================================================

\begin{itemize}
  \item \textbf{Y-14 is the calibration plate.} Measure it at the start and end
    of EVERY session. It is the anchor for all cross-session comparisons.

  \item \textbf{Arduino sensor must be on for ALL Helmholtz measurements.} This
    is the only way to get synchronous ambient temperature.

  \item \textbf{1337 = missing data.} If a measurement fails, record 1337 (not 0,
    not blank). The analysis scripts know to skip this value.

  \item \textbf{Record anomalies.} If something unusual happens (power glitch,
    sample dropped, instrument reset), note it in the session log with the time.
    Better to have too many notes than too few.

  \item \textbf{Do not mix sessions.} All 4 slots of a Y-plate must be measured
    in the same session. If you cannot finish a plate, restart it in the next
    session.

  \item \textbf{Save data immediately.} After each plate, verify the data file
    was saved. Data loss from a 9-hour session would be devastating.
\end{itemize}


\end{document}
